\subsection{Construction of the Feature Catalogue}
\label{subsec:feature-catalogue}

The functional scope defined in
Section~\ref{subsec:functional-scope} identifies the categories considered in
the evaluation. However, these categories remain too broad to support a
feature-level comparison between GEAR and the selected frameworks. We
therefore constructed a normalized feature catalogue that decomposes each
retained category into concrete and independently assessable features. This
catalogue constitutes the common reference used to build the coverage
matrices.

\subsubsection{Feature Identification}

Features were identified from two complementary sources. For GEAR, we examined
the metamodel, the feature models, the connector implementations, and the
generated code. For each studied framework, we examined the versioned official
documentation, API references, and implementation examples associated with
the version reported in Section~\ref{subsec:framework-selection}. When the
documentation was insufficient to determine the semantics of a construct, its
implementation or an executable example was additionally inspected.

Let $\mathcal{F}_{G}$ denote the set of features identified in GEAR and
$\mathcal{F}_{j}$ the set identified in framework $j$. Before filtering and
normalization, the initial set of candidate features is defined as:

\begin{equation}
    \mathcal{F}_{\mathrm{raw}}
    =
    \mathcal{F}_{G}
    \cup
    \bigcup_{j=1}^{m}\mathcal{F}_{j},
    \qquad m = 10.
    \label{eq:raw-feature-catalogue}
\end{equation}

This union ensures that the catalogue is not restricted to the features
currently represented by GEAR. A feature exposed by only one framework is
therefore retained as a candidate when it belongs to the functional scope of
the study.

The extraction focused on explicit constructs that allow developers to
configure the agents, models, tasks, tools, memory, data, orchestration,
control flow, or execution properties of a multi-agent system. Each candidate
was recorded together with its original terminology and the documentation,
API element, source file, or executable example from which it was identified.

\subsubsection{Scope-Based Filtering}

The candidate features were subsequently filtered according to the functional
scope established independently in
Section~\ref{subsec:functional-scope}. Let $\mathcal{C}^{*}$ denote the set of
functional categories included in the study and
$\operatorname{cat}(f)$ the category assigned to a candidate feature $f$. The
scope-filtered set is defined as:

\begin{equation}
    \mathcal{F}_{\mathrm{scope}}
    =
    \left\{
        f \in \mathcal{F}_{\mathrm{raw}}
        \;\middle|\;
        \operatorname{cat}(f) \in \mathcal{C}^{*}
    \right\}.
    \label{eq:scope-filtered-features}
\end{equation}

Consequently, the catalogue excludes constructs that do not concern the
functional configuration, orchestration, or execution of the studied
multi-agent systems. These include:

\begin{itemize}
    \item provider-specific deployment services and proprietary cloud
    capabilities;
    \item billing, account-management, and authentication mechanisms;
    \item framework-specific graphical interfaces;
    \item observability, monitoring, and execution-analytics services;
    \item security and governance mechanisms; and
    \item general-purpose data-processing capabilities that are not directly
    involved in the design or orchestration of agents.
\end{itemize}

These exclusions do not imply that the corresponding capabilities are
unimportant. They indicate that they belong to dimensions that require
different evaluation objectives and protocols. In particular, observability
is addressed separately in the operational evaluation and is therefore not
included in the functional coverage catalogue.

\subsubsection{Semantic Normalization}

The terminology used by the frameworks cannot be compared directly. Different
terms may refer to equivalent capabilities, while similar terms may conceal
different execution semantics. We therefore normalized the extracted features
according to their meaning rather than their names.

Two candidate features were merged when they represented the same
configuration decision and produced equivalent observable effects on the
structure or execution of the multi-agent system. Conversely, they were kept
separate when their differences affected at least one of the following
properties:

\begin{itemize}
    \item the entities involved, such as agents, tasks, tools, or workflow
    nodes;
    \item the direction and propagation of data;
    \item the ordering or concurrency of execution;
    \item the conditions governing branching, routing, or termination;
    \item the creation, persistence, or sharing of state;
    \item the aggregation of intermediate results; or
    \item execution constraints such as iteration limits, timeouts, retries,
    or human intervention.
\end{itemize}

For example, sequential execution and task dependency were retained as
distinct features. Sequential execution specifies an execution order, whereas
a dependency additionally expresses that one component relies on the
completion or result of another. Similarly, individual memory, shared memory,
and workflow state were not merged because they differ in ownership,
visibility, and lifetime.

Let $\sim_{\mathrm{sem}}$ denote the semantic equivalence relation established
through these criteria. The final normalized catalogue is the quotient set:

\begin{equation}
    \mathcal{F}^{*}
    =
    \mathcal{F}_{\mathrm{scope}}
    \big/
    {\sim_{\mathrm{sem}}}.
    \label{eq:normalized-feature-catalogue}
\end{equation}

Each equivalence class in $\mathcal{F}^{*}$ represents one normalized feature,
independently of the terminology used by GEAR or the individual frameworks.
This normalization avoids counting lexical variants as separate features
while preserving distinctions that may affect coverage or semantic fidelity.

\subsubsection{Catalogue Structure and Traceability}

Each normalized feature is documented through:

\begin{itemize}
    \item a unique identifier;
    \item its functional category;
    \item a framework-independent name and definition;
    \item the semantic criteria that must be satisfied for the feature to be
    considered covered;
    \item the technologies in which the feature was identified; and
    \item the corresponding evidence, including documentation references, API
    elements, source files, examples, or executable tests.
\end{itemize}

Table~\ref{tab:functional-catalogue-example} provides an illustrative excerpt
from this structure. The complete catalogue, including the provenance and
evidence associated with every feature, is provided in the replication
package.

\begin{table*}[t]
    \centering
    \caption{Illustrative structure of the normalized feature catalogue.}
    \label{tab:functional-catalogue-example}
    \footnotesize
    \begin{tabularx}{\textwidth}{
        @{}
        l
        p{0.15\textwidth}
        p{0.16\textwidth}
        >{\raggedright\arraybackslash}X
        p{0.27\textwidth}
        @{}
    }
        \toprule
        \textbf{ID}
        & \textbf{Category}
        & \textbf{Feature}
        & \textbf{Definition}
        & \textbf{Semantic criteria} \\
        \midrule

        AG1
        & Agent specification
        & Agent instructions
        & Definition of the behaviour expected from an agent
        & Instructions are explicitly associated with a specific agent and
        influence its behaviour. \\

        TK1
        & Tasks and planning
        & Task dependency
        & Requirement that one task be completed before another task can
        proceed
        & The dependency constrains execution and, when required, propagates
        the result of the preceding task. \\

        OR1
        & Orchestration and topology
        & Sequential execution
        & Execution of agents, tasks, or nodes according to a defined order
        & The order is preserved and a component starts only after its
        predecessor has completed. \\

        OR2
        & Orchestration and topology
        & Parallel execution
        & Concurrent execution of independent agents, tasks, or branches
        & Independent components may execute concurrently without an imposed
        sequential order. \\

        OR3
        & Orchestration and topology
        & Fan-in
        & Synchronization of results produced by multiple execution branches
        & Downstream execution waits for and combines the required upstream
        results. \\

        CT1
        & Control flow
        & Bounded loop
        & Repetition of an execution block subject to a maximum number of
        iterations
        & The loop may repeat a component while enforcing an explicit upper
        bound. \\

        \bottomrule
    \end{tabularx}
\end{table*}

The resulting catalogue defines the unit of analysis used throughout the
functional coverage evaluation. It supports the construction of two
complementary matrices: one determining which GEAR features can be generated
for each target framework, and another determining which framework features
can be represented by the GEAR metamodel. The construction and interpretation
of these matrices are presented in the following subsection.

